\pdfoutput=1
\documentclass[11pt]{article}
\usepackage[a4paper,margin=23mm,headheight=15pt]{geometry}
\usepackage[T1]{fontenc}
\usepackage{lmodern,microtype,amsmath,amssymb,amsthm,booktabs,array,longtable,enumitem,fancyhdr,xurl,needspace,xcolor}
\usepackage[colorlinks=true,linkcolor=blue!40!black,urlcolor=blue!40!black,citecolor=blue!40!black]{hyperref}
\hypersetup{pdftitle={ST215: direct transfer through contextual reassociation},pdfauthor={Research contribution for the ProveIt project}}
\setlength{\parindent}{0pt}
\setlength{\parskip}{5pt}
\setlength{\emergencystretch}{3em}
\setlist{nosep,leftmargin=1.6em,topsep=3pt}
\pagestyle{fancy}\fancyhf{}
\fancyhead[L]{\small ST215: direct transfer through contextual reassociation}
\fancyhead[R]{\small 21 October 2026}
\fancyfoot[C]{\thepage}
\newtheorem{theorem}{Theorem}
\newtheorem{lemma}[theorem]{Lemma}
\newtheorem{proposition}[theorem]{Proposition}
\newcommand{\pair}{\mathsf P}
\newcommand{\imp}{\mathsf I}
\newcommand{\quant}{\mathsf Q}
\newcommand{\AX}{\mathsf{AX}}
\newcommand{\MP}{\mathsf{MP}}
\newcommand{\Tr}{\mathsf S}
\newcommand{\calD}{\mathcal D}
\newcommand{\Con}{\operatorname{Con}}
\newcommand{\Halt}{\operatorname{Halt}}
\newcommand{\file}[1]{{\small\path{#1}}}
\newcommand{\hashline}[1]{\begin{center}\footnotesize\path{#1}\end{center}}
\begin{document}
{\LARGE\bfseries ST215: direct transfer through contextual reassociation}\par
\vspace{6pt}
{\large A materially different 215-state consistency machine}\par
\vspace{4pt}
Research contribution for Vladimir Reshetnikov's ProveIt project\\
Prepared with AI-assisted construction and independent computational and mathematical review\\
21 October 2026

\begin{abstract}
We give a deterministic, one-tape, two-symbol machine with 215 working
states whose blank-tape nonhalting is equivalent to $\Con(\mathrm{ZFC})$,
by an ordinary proof formalizable in ZFC and using the repository's
established numerical-calculus calibration. The new implementation
realizes a direct, 35-leaf transfer command previously blocked by two
physical source-context constraints. Reassociating the outer source
tree satisfies both constraints without changing any edge of its
complete ordinal control-flow graph. The source has 201 distinct
subtrees and 4,732 primitive occurrences. All 6,151 universal whole-tape
primitive cases pass; every one of the 430 literal transition rows is
covered. An exact 70,506-step blank run matches the interior of an
axiom command after a 177-step source prefix. Completing that command
places ten theorem heads above an opaque stack tail; the scheduler
requires only two heads of protected credit. The new table is distinct,
up to rooted state relabeling and reflection, from the compared
215-state tables in the pinned repository and twelve incoming archives.
The 214-state target is not reached. We also give a command-language
obstruction for rotating pop, correct a false return edge in an earlier
scratch-biased experiment, and prove an abstract adjacent-candidate
synchronization lemma that may support future encodings. The archive
contains exact tables, editable TeX, the PDF, proofs, deterministic
reconstruction, complete startup data, and offline certificate replay.
\end{abstract}
\setcounter{tocdepth}{1}
\begingroup\footnotesize\setlength{\parskip}{0pt}\tableofcontents\endgroup
\newpage

\section{Result, conventions, and novelty}
\label{sec:result}

The deliverable is the new machine \textbf{ST215}. It has 215 working
states and a separate halting label, with 430 total transition rows.
The native table starts in state $0$, uses dispatcher $84$, and halts
on entering $254$. The dense presentation uses working states
$0,\ldots,214$, start $0$, and halt $215$. Both start with the head at
cell $0$ of a bi-infinite all-zero tape; every transition writes one
bit and moves one cell left or right. No initial program or data are
supplied externally.

\begin{theorem}[Main result]
\label{thm:main}
Using the pinned inherited numerical-calculus calibration and
all-natural arithmetic macro theorems, the construction and proofs
in this packet establish
\[
 \mathrm{ZFC}\vdash
 \bigl(\neg\Halt(\mathrm{ST215},\mathrm{blank})
       \longleftrightarrow \Con(\mathrm{ZFC})\bigr).
\]
\end{theorem}

This is an ordinary mathematical result with replayable symbolic and
finite certificates. We supply no newly compiled exact-table Lean or
Rocq endpoint. The published K215 proof has the same ordinary-versus-kernel
distinction at the pinned snapshot \cite{repo,k215}. The foundational
calibration is reused, not replaced by a newly invented consistency
assumption.

\begin{center}\small
\begin{tabular}{lrrr}
\toprule
Quantity & Published K215 & Composite parent & ST215\\
\midrule
Working states &215&215&215\\
Distinct source subtrees &202&201&201\\
Primitive occurrences &4,751&4,775&4,732\\
Guarded primitive cases &6,176&6,209&6,151\\
Transfer body occurrences &---&78&35\\
Symbolic inference steps &1,698,790&1,611,489&1,617,630\\
Minimum duration constant &6&7&5\\
Blank steps to certified frame &95,714&68,442&70,506\\
\bottomrule
\end{tabular}
\end{center}

The transfer parent is the incoming T215 packet's frozen
\file{reference/transfer-parent}, not the published K215 machine
\cite{transfer}. Smaller source size does not automatically mean fewer
physical states: tree dispatch, arithmetic helpers, register scans,
and initialization share states under context-sensitive contracts.
Nor does the smaller primitive count prove faster execution.

The material change has three parts. First, the command language has
only transfer, modus ponens, and axiom generation, with field arities
$1,2,3$. Second, transfer is now the direct 35-leaf implementation
testing only the incoming cache, in place of the 78-leaf composite
testing both top and cache. Third, the source is placed in a new
outer tree context, with a zero-valued spare-block refill and a
different calibrated blank startup. It is not a state renaming of an
existing compared table.

The native table's SHA-256 is
\hashline{d5237504325eee96b0040d4d9574731a1caf4bcda3d6a49e94827d9532a664e1}
The dense table's SHA-256 is
\hashline{ead0ec3aa1c3502015faa8a077ded98be7522fa907a9df43a8052de555398375}
The source's SHA-256 is
\hashline{212cf0ffde43c1cc22c183dd8bf0b3a699396f96c9942ee92a7cc372cd273ca9}

\section{Numerical foundation: no formula-decoding requirement}
\label{sec:foundation}

For natural numbers put
\[
 \pair(a,b)=\frac{(a+b)(a+b+1)}2+b,\qquad
 \imp(a,b)=4\pair(a,b),\qquad
 \quant_a b=4\pair(a,b)+2.
\]
Write $a\to b=\imp(a,b)$, falsity $F=1$, truth $0$, and
$NX=\imp(X,1)$. Also put
$M(a,b)=4\pair(a,b)+1$, $E(a,b)=4\pair(a,b)+3$, and
$\exists_a X=N\quant_a NX$.
Parameters range over \emph{all natural numbers}. No formula parser,
variable-support test, or well-formedness filter is inserted. In
particular $\imp(0,0)=0$, and $\operatorname{unpair}(0)=(0,0)$.
Those aliases are used as literal numerical equalities.

The retained calculus $\calD$ has truth, numerical MP, theorem-only
GEN, seven logical rows, the original gated bundle, and the dominated
row below:
\begin{align*}
 R_1(a,b)&=\exists_aE(a,b),\\
 R_2(a,b,c)&=\quant_a(b\to c)\to(\quant_ab\to\quant_ac),\\
 R_3(a,b,c)&=E(a,b)\to(M(c,a)\to M(c,b)),\\
 R_4(a,b,c)&=\quant_a\quant_bc\to\quant_b\quant_ac,\\
 R_5(a,b,c)&=E(c,a)\to(M(c,b)\to M(a,b)),\\
 R_6(a,b,c)&=(\exists_a\quant_ac)\to c,\\
 R_7(a,b,c)&=E(c,a)\to(E(c,b)\to\quant_{M(a,b)}E(a,b)),\\
 L(A,B,C,U)&=((C\to U)\to A)\to((A\to C)\to(B\to C)),
 \qquad U\le A.
\end{align*}
The unused $b$ in $R_6$ is intentional and agrees with the literal source.
The enabled bundle is retained in full:
\begin{align*}
 B_0(b,k,m)&=\quant_b(m\to N\quant_b(m\to NM(b,k))),\\
 S_0(a,b,k)&=\quant_k(M(k,a)\to M(k,b)),\\
 C_0(a,b,k,m)&=\quant_b\bigl(M(b,k)\to
       N\quant_a(m\to N\quant_a(S_0(a,b,k)\to B_0(b,k,m)))\bigr),\\
 J_0(a,b,c)&=\quant_c(C_0(a,b,c,c)\to NB_0(b,c,c)),\\
 V_0(a,b,c)&=M(a,b)\to\quant_cM(a,b),\\
 \operatorname{Bundle}(a,b,c)&=N(V_0(a,b,c)\to J_0(a,b,c)).
\end{align*}
It is an axiom when $c\ne a,b$. The gate outputs truth otherwise.
Removing a support-bearing part of this bundle would require a new
calibration; we do not do so.

The inherited finite recovery certificate proves that the dominated
calculus has the same natural-code theorem set as its unrestricted
predecessor. It contains 3,728 proof nodes and 5,160 shared term nodes.
The included checker verifies explicit substitutions, MP, bijective
renamings, and exact endpoints. A dominated leaf is justified because
its $U$ is zero or a literal subterm of $A$; numerical implication
dominates its arguments, so every natural substitution preserves
$U\le A$. The only term normalization is $\imp(0,0)=0$.
Universal schema instantiation in this finite metaproof does not
silently add a substitution inference rule to the numerical machine.

The further inherited calibration identifies numerical refutation with
ZFC refutation:
\begin{equation}
 \calD\vdash1
 \quad\longleftrightarrow\quad
 \exists p\;\operatorname{Proof}_{\mathrm{ZFC}}(p,\bot).
 \label{eq:calibration}
\end{equation}
This is the established foundation of the repository's ordinary K215
proof \cite{k215,dominated}. It is stated here as a reused theorem,
not re-proved from first-order syntax. The distinction matters:
checking the 3,728-node propositional certificate alone does not
establish the whole comparison with ZFC.

Thus the user-permitted freedom to use junk codes is already present
in the actual formalism. The advance in this packet comes from a
smaller operational presentation and its physical implementation,
rather than from a newly claimed axiomatization of set theory.

\section{Three commands and the direct transfer}
\label{sec:commands}

Write a command store as $(t,w,d,h)$: top, encoded stack, additive
counter, and cached formula. The notation $[x;w]=\pair(x,w)$ extends
by right nesting to a finite explicit list over any numerical tail.
For $\operatorname{unpair}(w)=(x,v)$ the new command is
\begin{equation}
 \Tr(t,w,d,h)=
 \begin{cases}
 \mathsf{halt},&h=1,\\
 (h,v,d,x),&h\ne1.
 \end{cases}
 \label{eq:directS}
\end{equation}
An incoming top $t=1$ alone does not halt. The predecessor composite
$\mathrm{DROP};\mathrm{NIP};\mathrm{NIP}$ had the additional guard
$t=1$. Both forms have the same nonhalting store map, but the
single-guard contract is the exact one for which the original
three-command strong compiler was proved.

For $y=d+x$, modus ponens is
\[
 \MP(x)(t,w,d,h)=
 \begin{cases}
 (y,[0;w],y,0),&\imp(t,y)=h,\\
 (0,w,y,0),&\imp(t,y)\ne h.
 \end{cases}
\]
For $\AX(c;x,y)$ put $U=d$, $A=d+x$, $B=h+y$, $C=c$. Its output is
\begin{equation}
 \bigl(G(A,B,c),[R_1,\ldots,R_7,L(A,B,c,U),\quant_Uh,t;w],0,0\bigr),
 \label{eq:AX}
\end{equation}
where $R_j=R_j(A,B,c)$ and $G$ is the full bundle or zero according
to its gate. Here $h$ in $\quant_Uh$ is the incoming cache, before
adding $y$. The seven rows, $L$, and the gate are implication-coded,
and hence not $1$; $\quant_Uh$ has tag $2$ and is neither $0$ nor $1$.
The tenth head saves the incoming top.

The field lists are exactly
\[
 f(\Tr)=[0],\qquad f(\MP(x))=[1,x],\qquad
 f(\AX(c;x,y))=[c+2,y,x].
\]
The source's logical registers, in order, are
$(N,T,W,A,D,H,X,Y,L)$, with $N$ the next candidate, $A$ the
selector, $X,Y$ scratches, and $L$ the residual command pointer.
At a complete command boundary $A=X=Y=0$.

\paragraph{Why the direct body is only 35 occurrences.}
At source ordinal 4697, clear $T$ in a two-leaf loop. Next, while
decrementing $H$, increment both $T$ and $X$; move $X$ back into
$H$. This seven-leaf copy leaves $T=H=h$ and $X=0$. Apply the
unchanged 26-leaf NIP. That body tests $H=1$; otherwise it clears
$H$, unpairs $W$, leaves the tail in $W$, and loads the head into
$H$. Thus the total is $2+7+26=35$, and the map is
\eqref{eq:directS}. All copying loops decrease a natural register.
The inherited NIP theorem handles its guard and unpairing on every
natural input, including $W=0$.

The direct body preserves $N,L,D$ and returns zero $A,X,Y$ on a
nonhalting exit. Its only halt is primitive CH at ordinal 4717.
The replacement saves 43 occurrences relative to the 78-leaf
composite, while maintaining 201 distinct source subtrees.

\section{The contextual reassociation that makes it fit}
\label{sec:rotation}

The incoming short-transfer experiments already identified the
35-leaf semantic body. A straightforward placement failed its
physical CD context; a subsequent context-root placement fixed
CD but failed its return continuation. The present construction
repairs the \emph{outer source context}, without an additional
physical helper or a modified candidate scheduler.

Write the context-root tree as
\[
 t=(A,((B,C),D)),
\]
where $A=(\mathrm{scheduler},\mathrm{field\ reader})$, $B$ contains
the selector and AX/MP bodies, $C$ clears top, and $D$ copies cache
then performs NIP. Set
\begin{equation}
                  t'=(((A,B),C),D).
 \label{eq:rotation}
\end{equation}
This is not an arbitrary associativity law for a structured-control
language: ancestor-relative breaks and continues can change their
targets when the tree changes. Here preservation is checked for
\emph{every} guarded source edge by two independently implemented
resolvers.

\begin{center}\small
\begin{tabular}{llll}
\toprule
Purpose & Ordinal & Old address & New address\\
\midrule
Return to scheduler, CC3 &4696&10011&00111\\
Clear top, CD1 &4697&1010&010\\
Unique halt, CH &4717&\multicolumn{2}{l}{new address 11101}\\
\bottomrule
\end{tabular}
\end{center}

CC3 still chooses its third right ancestor, which is the root.
Its new address ends in $111$, the suffix required by the existing
compressed physical continuation helper. CD1 has the required
preceding marker bit in its new context. These local observations
explain the repair; the universal whole-tape runtime certificate is
what establishes the physical implementation.

All 4,732 source addresses change relative to context-root, while
the ordered opcode/operand list is identical. The two resolvers agree
on all 6,151 guarded cases. The complete ordinal CFG also equals
that of the incoming simple direct-transfer source. Only the
syntactic spelling of its return differs there; the target is
unchanged. The tree is closed, has 201 distinct subtrees, 4,731
internal prefix occurrences, maximum depth 44, and maximum zero
run 6.

The deterministic reconstruction begins with the frozen transfer
parent source, extracts its seven-leaf cache copy and 26-leaf NIP,
forms the direct body, makes the context-root transformation, and
then applies \eqref{eq:rotation}. The native table is reconstructed
from the parent's literal table by an explicit row patch. The
noteworthy final row
\[
                         198\quad0\quad0\quad L\quad40
\]
makes newly appended spare blocks zero-valued. The source semantics
are independent of spare values, but initialization is not; the
blank proof below concerns this exact row and start state $0$.
The full recipe is executable in \file{checks/rebuild.py}.

\section{Whole-tape physical simulation}
\label{sec:runtime}

The certificate covers every primitive occurrence, with separate
zero and positive cases for each decrement. It is not a collection
of bounded integer tests. Each case quantifies over all natural
primary register values, all nonnegative gaps, and every nonempty
finite list of spare values.

A standard frame, with dispatcher head at position $p$, has a
left marker at $p-1$, the finite source-address bits beginning at
$p$, a zero gap to a header marker at $p+64+g$, and the register
file immediately to its right. Register blocks are
\[
             B(v)=0\,1^{v+1}.
\]
Logical registers map to physical block indices by
\[
                    (6,2,0,4,3,7,8,5,1).
\]
After the nine primary blocks comes a nonempty arbitrary finite
spare list; all remaining tape is zero. The certificate proves
the complete tape relation, not only selected registers.

The proof engine represents run lengths by affine expressions with
13 coordinates: a constant coordinate, gap, nine primary values,
first-spare value, and the width of the remaining spare list.
All variable coordinates denote nonnegative values.
It uses three sound steps:

\begin{enumerate}
\item A literal step checks a row of the exact table and rewrites one
cell, with exact head motion and elapsed time.
\item A unary step proves a directional self-loop over a homogeneous
run by induction on its arbitrary natural length.
\item A list step applies a separately derived one-block schema by
structural induction over an arbitrary finite list of blocks.
\end{enumerate}

The four one-block schemas have duration $v+2$. They scan $B$
right, convert $B$ to $C=1\,0\,1^v$ right, scan $B$ left, or convert
$C$ back to $B$ left. The list normalization
$C(A)\,1=1\,B(A)$ follows by induction on the block list $A$.
Opaque left and right contexts in the one-block proofs ensure
that no unexamined tape cell is changed.

Every symbolic case is derived from the literal table after an
explicit injective relabeling to the engine's historical labels.
An increment appends a zero spare. No primitive deletes
the final spare; the represented frame class is therefore closed.
The output address and gap changes agree with source successor,
break, continue, increment, and decrement semantics. The unique
CH case reaches the external halting label.

\begin{center}
\begin{tabular}{lr}
\toprule
Exact-table obligation & Observed result\\
\midrule
Guarded cases proved &6,151 of 6,151\\
Literal table rows covered &430 of 430\\
Symbolic inference steps &1,617,630\\
Smallest duration constant &5\\
Failed cases &0\\
Independent block schemas &4\\
\bottomrule
\end{tabular}
\end{center}

Each duration has nonnegative coefficients and a positive constant.
Thus every source step takes a finite positive number of binary
steps, including the halting case. Infinite source execution
cannot collapse to a finite physical execution. A finite complete
source run lifts to a finite binary run.

The saved receipts are \file{evidence/runtime.json} and
\file{evidence/runtime_cases.json.gz}. The default replay audits
their bindings, order, coverage, durations, and four block schemas.
The explicit \texttt{--runtime} option rederives all 6,151 cases
and compares every resulting case record with the saved record.
Python and the ordinary soundness argument are part of this
verification boundary; these are not proof-assistant kernel
certificates.

\section{Exact blank startup and the partial AX bridge}
\label{sec:startup}

Literal execution from state $0$ on blank tape reaches the following
complete frame after exactly \textbf{70,506 binary transitions}:
\[
\begin{gathered}
 q=84,\qquad p=-2,\qquad \mathrm{header}=62,\qquad g=0,\\
 \mathrm{PC}=00110010010110010,\qquad \mathrm{ordinal}=58,\\
 (N,T,W,A,D,H,X,Y,L)=(1,0,0,0,0,0,0,18,0).
\end{gathered}
\]
There are 63 spare blocks, every one zero-valued, and exactly 100
one-cells on the entire tape. The full one-cell set is recorded in
\file{evidence/bootstrap.json}; every other cell is zero. The
independent native and dense executions agree at every step.
\file{evidence/blank-delta.jsonl.gz} contains all 70,506 transition
deltas, not merely sampled snapshots.

This is an interior AX frame, not a clean source-root configuration.
To identify it, execute the source from AX entry ordinal 25 with
\[
                (N,T,W,A,D,H,X,Y,L)=(1,0,2,0,0,0,0,0,0).
\]
After exactly 177 primitives it reaches ordinal 58 and the complete
displayed logical store. Both CFG resolvers produce the same trace.
The full before/after stores, literal instruction, and next ordinal
are stored for every source transition in
\file{evidence/AX-prefix.json}. An additional independent checker
constructs the physical frame directly from the native table and
the register-block definition.

The inherited all-natural AX theorem now applies to the
\emph{remaining suffix} of this exact command. Its synthetic input
is $\AX(0;0,0)$ on $(t,w,d,h)=(0,2,0,0)$. Completion terminates
and yields
\[
\begin{gathered}
 N=1,\quad L=0,\quad A=X=Y=0,\quad t=d=h=0,\\
 w=[R_1(0,0,0),\ldots,R_7(0,0,0),
       L(0,0,0,0),\quant_0 0,0;\,2].
\end{gathered}
\]
All ten explicit heads are native theorems: seven retained axiom
instances, the dominated row, a generalization of truth, and truth.
The gate is zero because $c=A=B=0$.

The suffix $2$ is deliberately treated as opaque. Indeed
$2=\pair(0,1)$ and $1=\pair(1,0)$, so blindly popping it would
eventually expose falsity. The credit theorem in
Section~\ref{sec:credit} prevents this. Merely reaching a
plausible register frame, or calling the tail harmless, would not
establish soundness.

It is unnecessary to literally simulate the enormous formula codes
generated by the rest of AX. Exact matching of a finite prefix,
the all-natural terminating AX theorem, and positive-duration
primitive simulation determine its future completion. No
asymptotic or probabilistic startup assumption is used.

\section{A complete strong compiler for the direct command}
\label{sec:compiler}

For completeness, the compiler must work over \emph{arbitrary}
retained top and stack, not only theorems. For every
$\calD$-theorem $X$, construct a finite command block $B_X$ such that
\[
 (t,w,0,0)\xrightarrow{B_X}\mathsf{halt}
 \quad\text{or}\quad (X,w,0,0)
 \qquad\text{for all natural }t,w.
 \qquad\text{(Strong)}
\]
Early halting counts as success. Soundness of actual runs is a
separate invariant. The inherited compiler \cite{transfer} was
proved for exactly the single-guard contract \eqref{eq:directS};
there is no need to infer it backwards from the more eager
composite guard.

\subsection{Structural macros}

Put $D=\AX(0;0,0)$ and
\[
 H=D;\Tr^{11},\qquad
 \operatorname{PUSH}=D;\Tr^9;\MP(0).
\]
After ten transfers following $D$, the tuple is
$(\quant_d h,w,0,t)$. None of the first ten tests sees a new
falsity head: only zero and the first nine generated entries have
been tested. The eleventh transfer halts precisely if $t=1$;
otherwise, writing $w=[x;v]$, it yields $(t,v,0,x)$.
Thus $H$ preserves top, pops one preexisting head, resets the
counter, and tests the \emph{incoming top}. It is not a substitute
for the old NIP's incoming-cache test.

After nine transfers in PUSH the tuple is
$(L,[t;w],0,\quant_d h)$. Its final MP necessarily fails because
the cache has tag $2$ whereas $\imp(L,0)$ has tag $0$. Hence
\[
             \operatorname{PUSH}(t,w,d,h)=(0,[t;w],0,0)
\]
and PUSH never halts, even when the original $t$ or $h$ equals $1$.

Define
\[
 \operatorname{DISCARD}
   =\MP(1);\AX(1;0,0);\Tr^{10};\MP(0).
\]
The first MP fails, since $\imp(t,1)>0$ while cache is zero;
it produces $(0,w,1,0)$. AX has $A=c=1$, gate zero, and saved
top zero. Ten transfers give $(\quant_1 0,w,0,0)$, after which
MP fails because $\imp(\quant_1 0,0)>0$. Thus DISCARD is $B_0$,
returning $(0,w,0,0)$ without halting.

When top, counter, and cache are zero, let $J(0)$ be the empty
word and $J(u)=\MP(u)$ for $u>0$. It sets the counter to $u$
without changing the stack. Omitting MP at $u=0$ is essential:
the alias $\imp(0,0)=0$ would make it succeed and push a zero.

\subsection{Axiom and generalization blocks}

For a retained row $R_j(a,b,c)$, $1\le j\le7$, use
\[
 \operatorname{DISCARD};\AX(c;a,b);\Tr^{j+1};H^{9-j}.
\]
The shifts put row $j$ on top. Each $H$ preserves that non-$1$
top and removes one generated stack entry; the last removed
entry is saved zero. The original stack is restored with zero
counter and cache.

For the dominated row $L(a,b,c,u)$ with $u\le a$, use
\[
 \operatorname{DISCARD};J(u);\AX(c;a-u,b);\Tr^9;H.
\]
For an enabled full bundle, use
\[
 \operatorname{DISCARD};\AX(c;a,b);H^{10}.
\]
The latter retains every part of the original gated bundle.

Given $B_X$, theorem generalization is compiled by
\[
 B_{\quant_iX}
   =B_X;\operatorname{PUSH};J(i);\Tr;\AX(i;0,0);\Tr^{10}.
\]
After PUSH and $J(i)$, the single transfer maps
$(0,[X;w],i,0)$ to $(0,w,i,X)$ without testing the newly loaded
$X$. AX has $U=A=c=i$, zero gate, ninth head $\quant_iX$,
and saved zero. Ten transfers return exactly
$(\quant_iX,w,0,0)$. This includes $i=0$ and $X=1$.

\subsection{Major-premise-first modus ponens}

For $M=\imp(A,B)$, define
\[
\begin{split}
 B_B={}&B_M;\operatorname{PUSH};B_A;\operatorname{PUSH};\\
       &\Tr;\Tr;\MP(B);\AX(B;0,0);\Tr^{11}.
\end{split}
\]
Unless a permitted earlier halt occurs, the first four blocks
produce $(0,[A,M;w],0,0)$. Two transfers produce $(A,w,0,M)$;
if the second one instead sees $A=1$, it has already succeeded
by halting. MP now succeeds by exact numerical equality and
returns $(B,[0;w],B,0)$.

The following AX has gate zero and saves $B$. Ten transfers
leave $(\quant_B0,[0;w],0,B)$. The eleventh halts if $B=1$,
and otherwise returns $(B,w,0,0)$. When $A=B=0$, the major
is literally zero and MP deliberately succeeds; the final
transfer removes precisely the zero it pushed. No hidden
nonzero-code assumption enters the calculation.

Induct on a finite native proof, expanding a finite proof DAG
to a proof tree if necessary. Compiling the major first creates
no dependency cycle. The preceding cases prove (Strong).
In particular, $B_1;H$ halts from every $(t,w,0,0)$.

\subsection{Raw wire encoding and cofinal successful candidates}

Let $F(fs,r)$ right-fold a field list by Cantor pairing over
$r$, and put
\[
 \operatorname{tail}([])=0,\qquad
 \operatorname{tail}(k::L)=1+F(f(k),\operatorname{tail}(L)).
\]
A positive residual is decremented before its next command;
the first command of a fresh candidate is read directly from
the candidate number. The successor outside $F$ is therefore
required for a continuation, but absent at fresh fetch.

Let $R=\operatorname{tail}(B_1;H)$. For every natural $Y$, the
candidate
\begin{equation}
                  N_Y=\pair(2,\pair(Y,\pair(1,R)))
 \label{eq:cofinal}
\end{equation}
starts with $\AX(0;1,Y)$, which resets counter and cache from
any retained store. The strong compiler accepts any resulting
top and stack, so the fixed suffix halts. Since $N_Y\ge Y$,
successful candidates are unbounded.

Each nonhalting positive residual strictly decreases: its next
value is a right projection after one, two, or three unpairings
of its predecessor. All arithmetic commands terminate on
natural inputs. Every nonhalting candidate therefore drains,
and ordinary candidate increment by one reaches any chosen
$N_Y$ above the initial cutoff. Completeness needs neither a
well-formedness filter nor a fairness hypothesis.

\section{Protected credit for the actual scheduler}
\label{sec:credit}

Here is the all-natural argument needed for the opaque startup
tail. Let $f$ denote the right Cantor projection and
$T_s=s(s+1)/2$. Let $W(n)$ be the entire syntactic command
word of fresh candidate $n$, ignoring possible semantic halts.
Set $R(0)=[]$ and $R(r)=W(r-1)$ for $r>0$. If
$\operatorname{unpair}(n)=(a,y)$, then
\begin{equation}
 W(n)=
 \begin{cases}
 \Tr;R(y),&a=0,\\
 \MP;R(f(y)),&a=1,\\
 \AX;R(f(f(y))),&a\ge2.
 \end{cases}
 \label{eq:schedule}
\end{equation}
The command parameters are irrelevant to credit, but their
arity and continuation decrement are not.

Assign weights $\Tr=-1$, $\MP=0$, $\AX=10$.
Successful MP's additional truth push is conservatively
ignored. Write $b(n)$ for total weight,
$B(k)=\sum_{n<k}b(n)$, and
$g(0)=0$, $g(r)=b(r-1)$ for $r>0$.
By induction $W(n)$ has at most $n+1$ commands, so every
prefix of $R(r)$ loses at most $r$ units.

\begin{lemma}[Cantor prefix pullback]
If a sequence has nonnegative initial partial sums, its
composition with $f$, and then with $f^2$, has nonnegative
initial partial sums.
\end{lemma}
\begin{proof}
In increasing argument order, $f$ concatenates the lists
$[0],[0,1],[0,1,2],\ldots$. Every finite prefix consists of
full lists followed by an initial partial list. Each
corresponding sum is nonnegative. Repeat for $f^2$.
\end{proof}

\begin{lemma}[Later diagonals require no new credit]
Every Cantor diagonal $s\ge2$, read in increasing candidate
order, has nonnegative weight at every command prefix,
relative to its start. Moreover $B(k)+k\ge0$ for every $k$,
and $B(T_s)\ge-3$.
\end{lemma}
\begin{proof}
The exceptional words are $W(0)=\Tr$, $W(1)=\MP$,
and $W(2)=\Tr\Tr$. The first values of $B$ are therefore
$0,-1,-1,-3$, supplying the induction base.

Use the shifted sequence $q(r)=g(r)+1$. For $k>0$,
\[
                   \sum_{r<k}q(r)=k+B(k-1)\ge0.
\]
On diagonal $s\ge2$, the first $s-1$ candidates begin
with AX, the next begins with MP, and the last with S.
The first $k$ AX candidates have total weight at least
$9k$ by the twice-pulled-back inequality. Before AX
at position $y$, the balance is at least $9y$; after
its initial $+10$, any suffix prefix loses at most
$f^2(y)\le y$. Thus all AX prefixes are safe.

After that portion, balance is at least $9(s-1)$.
The MP suffix loses at most $s-1$, and the last S
and its suffix at most $1+s$. Hence every prefix
through those last two words is bounded below by
\[
                    9(s-1)-(s-1)-1-s=7s-9>0.
\]
All pullback indices used here belong to earlier
diagonals, so the induction is well founded.
The new diagonal is prefix-safe; combining
$B(T_s)\ge-3$ with $k\ge3$ gives $B(k)+k\ge0$ at
its new boundaries. This closes the joint induction.
\end{proof}

Starting at candidate $1$, the exceptional prefix is
$W(1)W(2)=\MP\Tr\Tr$, with minimum weight $-2$.
No later diagonal makes this worse. Thus the exact
\emph{conservative} debt at the actual cutoff is
\[
                             d(1)=2.
\]
This is a statement about the syntactic weight bound;
it does not claim every actual execution attains
that debt, since successful MP may add heads.

Maintain an explicit finite list of theorem heads
above a fixed opaque tail $J$. AX prepends ten heads,
MP prepends a theorem zero on success and none on
failure, and S consumes one head. With two initial
heads the list has enough credit before every S.
Top and cache remain theorem-valued; the S guard
can therefore halt only on a native theorem $1$.
An explicit list is needed because $\pair(0,0)=0$
does not uniquely encode list length. The invariant
uses only
$\operatorname{unpair}(\pair(x,v))=(x,v)$.

ST215 begins the ordinary schedule with ten heads,
so at least eight explicit heads remain over $J=2$
at every boundary. Its unsafe numerical tail never
becomes visible. The full inherited proof is
included as \file{proofs/protected_credit.md}.

\section{The consistency equivalence}
\label{sec:equivalence}

\begin{proof}[Proof of Theorem~\ref{thm:main}]
The exact finite blank trace reaches the represented
interior AX frame of Section~\ref{sec:startup}.
The matching source prefix and all-natural AX
contract complete it to ten theorem heads above
tail $2$, with next candidate $1$ and empty
residual. The universal physical simulation
transfers this conclusion to the literal table.

For soundness, the credit theorem protects the
tail, and the command equations preserve
theorem-valued top, cache, and explicit heads.
The only halt tests cache equal to $1$.
Thus binary halting gives $\calD\vdash1$, and
\eqref{eq:calibration} gives a finite ZFC refutation.

Conversely, a ZFC refutation gives $\calD\vdash1$
by the same inherited calibration. The strong
compiler produces $B_1;H$. Equation
\eqref{eq:cofinal} supplies arbitrarily large
successful candidate numbers, irrespective of
the retained stack. The ordinary exhaustive
scheduler reaches one unless it halts earlier.
Source termination of that successful candidate
lifts to finite binary halting.

These are finite calculations, primitive-recursive
proof and word transformations, and induction on
natural numbers, finite lists, and derivations.
They are formalizable in ZFC without assuming
its consistency. Negating the resulting halting
equivalence yields the displayed nonhalting
equivalence.
\end{proof}

In particular, this result is not an unconditional
proof that ST215 never halts. Assuming the usual
consistency of ZFC, it does not halt and ZFC
cannot prove that fact. No claim about ZFC's
inability to prove a false halting statement
from mere consistency is needed here.

\section{Other formalisms and research results}
\label{sec:research}

The search was pursued at two levels: changing the operational
calculus and changing its physical realization. The positive
machine uses the inherited calibrated numerical foundation,
but dispenses with a separate DROP command. The following
results explain why some more radical small interfaces fail
and identify a different abstract scheduler.

\subsection{Rotating pop has a structural obstruction}

One incoming experiment replaces transfer by
\[
 \mathsf R(t,[x;v],d,h)=(x,v,d,t),
\]
halting instead when $h=1$ or $x=1$. A compact runtime does
not make this interface refutationally complete.

Call an adjacent pair $(m,a)$ productive when
$m=\imp(a,b)$ for some $b>0$. Consider the spine consisting
of cache, top, then successive stack heads. A spine is
blocked if it contains neither $1$ nor a productive pair.
The all-zero spine is blocked. R drops its first spine
entry. An MP success on a blocked spine must have
consequent zero. The literal AX row order creates no
productive adjacency: the ten necessary mismatches
follow from injectivity and incompatible outer tags
modulo four. Thus every command word from a zero store
remains blocked and never halts.

This excludes the proposed strong compiler even for the
fixed native theorem
\[
                    20=\imp(0,\quant_0 0)=\imp(0,2).
\]
Truth, GEN, the recovered K schema, and MP prove it.
A returned store $(20,0,0,0)$ would contain the productive
adjacent pair $(20,0)$, while the alternative successful
outcome, halting, is also excluded. This obstruction is
unconditional and does not assume $\Con(\mathrm{ZFC})$.
Appendix~\ref{sec:rotating-obstruction} supplies the full
all-natural proof and its precise scope for the archived
physical experiment.

\subsection{A scratch increment changed the return target}

The incoming scratch-biased short-transfer experiment inserted
CI6 before CC3. Its earlier return address $10011$ became
$100111$. The extra right ancestor makes CC3 return to
selector ordinal 23 rather than scheduler ordinal 0.
This is a semantic change, despite an unchanged opcode.

After AX, selector and cache are zero. CI6 leaves scratch
$X=1$; the misdirected return immediately enters S.
The copy/restore loop injects that scratch unit into
cache, and NIP halts. The all-natural suffix length is
$2t+15$ source transitions; for the actual initial
zero gate it is 15. The exact source checks and proof
are in Appendix~\ref{sec:scratch-correction}.
The physical primitive checker can faithfully implement
this faulty source and still pass every runtime case.
It cannot prove the advertised higher-level parser.

\subsection{Consecutive successful candidates tolerate skipped numbers}

A different \emph{abstract} parser can be calibrated if its
equations are implemented correctly. Let a bit $b$ become
zero after S and one after MP or AX. Read the next command
from $L-1+b$ when the continuation pointer $L$ is positive.
Suppose fresh candidate numbers increase by one or two.
Unbounded successful candidates alone are insufficient:
a sparse search could skip all of them.

For any $r>0$, define
\[
 a=\pair(2,\pair(1,\pair(0,r))),\quad
 z=\pair(1,a),\quad n=\pair(1,z).
\]
Candidate $n$ runs
$\MP(1);\AX(0;0,1)$ and reaches $(L,b)=(r,1)$.
The identity
$\pair(1,z)+1=\pair(0,z+1)$ makes candidate $n+1$
run $\Tr;\MP(2);\MP(0)$ and reach the same parser
state, or halt sooner. The arithmetic stores need
not agree; the common suffix begins with a reset.

Encode a successful reset-plus-refutation suffix
with an arbitrary positive unused tail $q$.
This gives cofinally many \emph{consecutive}
successful pairs. Any increasing schedule with
gaps at most two meets each sufficiently late
pair. Positive continuation pointers strictly
decrease, so the search actually reaches its
subsequent selected candidates.

The positive tail is essential: with zero tail,
the biased encoding of a final S can itself be
zero, and the parser would skip it. The complete
arithmetic proof is Appendix~\ref{sec:biased-parser-synchronization}.

There is also a separate protected-credit theorem
for the hypothetical parser in
\file{research/biased_credit_proof.md}. Writing
$D(n)$ for a word's maximum debt, it proves
$2^{D(n)}\le n+2$ and
$\sum_{n<k}2^{D(n)}\le4k$ by Cantor-diagonal
induction. The pullback bound then makes every
selected diagonal $s\ge5$ prefix-safe for
\emph{any} selection with gaps one or two.
The explicit initial schedule from candidate
$1$, bit $1$, also has nonnegative balance.
With ten theorem heads this gives protected-tail
soundness. Combined with the consecutive-pair
compiler, these results calibrate the intended
abstract search equations.

No physical implementation of those biased equations
is claimed here. The archived scratch-biased machine
does not implement them, and ST215 does not use them.
Their potential value is that an implementation may
retain a scratch bit instead of paying to erase it,
while preserving the only required property:
equiconsistency. A source and table implementing
the equations remain a separate research problem.

\section{Physical comparison and the remaining 214-state problem}

The comparison corpus contains 78 distinct 215-state
\emph{byte files}, under 332 original file names,
drawn from the pinned K215 directory and twelve incoming
archives. Those are not claimed to be 78 mutually
nonisomorphic machines. Every input is included
verbatim with its origin and SHA-256.

A rooted deterministic labelled-graph traversal compares
write bits, directions, read-symbol edges, and halt
targets. It finds no state-relabeling isomorphism
between ST215 and any corpus table, including global
left--right reflection. All 215 working states are
reachable in this abstract graph. Graph reachability
does not assert blank-run reachability.

As a control, restore just the append-one write
at $198/0$. The resulting precursor still has no
direct or reflected match. Thus the difference is
not solely the refill bit. Its action differs from
published K215 after the abstract read word
\[
                   0000001001000000001001:
\]
the control writes $1$ and moves right where K215
writes $0$ and moves left. This read word is a
graph witness, not a claimed physically realized
tape history.

The canonical rooted graph hash of the final native
table is
\hashline{28eceb93cebfd2529c9140f0bf6c6d7c3bf110f921cf584958caa61efa5e02f4}
The corpus, algorithm, receipts, and independent
review are in \file{research/comparison/}.
This establishes novelty relative to the inspected
corpus, not global historical priority.

Bounded attempts to save the final state found no
accepted 214-state machine:
\begin{center}\small
\begin{tabular}{lrr}
\toprule
Search family & Specifications & Survivors\\
\midrule
K215 local folds, radius 5, fuel 97 &316,292&0\\
K215 local folds, radius 7, fuel 97 &316,292&0\\
Short-layout local folds, radius 7, fuel 97 &309,444&0\\
Concrete primitive row-supplier quotients &91,164&0\\
\bottomrule
\end{tabular}
\end{center}
The short-layout searches use the append-one
precursor, not the final append-zero table.
The quotient screen first passes all 6,151
baseline concrete primitive cases. Its finite
arithmetic fixtures, tape windows, and budgets
do not prove universal validity for any survivor,
and failure to find a survivor proves no lower
bound. Exact inputs, portable drivers, original
receipts, and limits are in
\file{research/physical/}.

The usable gain is a physically valid short
transfer source. Future state reduction can begin
with its 35-leaf body and calibrated startup.
The successful repair suggests searching over
source contexts jointly with state allocation:
a primitive's reusable physical implementation
can depend on ancestor bits that a source DAG
count alone does not record.

\section{Reproduction, artifacts, and inherited boundaries}
\label{sec:reproduction}

The packet is designed to be copied into the research
directory without requiring a repository checkout
to run its main checks. Python 3 and its standard
library suffice. Assertions must be enabled.

\begin{verbatim}
python3 -I verify.py --work /tmp/st215-finite
python3 -I verify.py --runtime --work /tmp/st215-full
\end{verbatim}
The selected output directory must not already exist.
The first command verifies the manifest,
deterministic reconstruction, both CFG resolvers,
source identities, exact native/dense startup,
the source AX prefix, the fixed propositional
certificate, saved runtime metadata, and the
independent graph comparison. The second also
rederives every symbolic primitive case. On the
preparation machine the latter symbolic derivation
takes a few minutes; timing is not a theorem.

\begin{center}\small
\begin{tabular}{p{.28\linewidth}p{.64\linewidth}}
\toprule
Path & Purpose\\
\midrule
\file{machine/} & Exact native/dense ST215 tables and explicit bijection.\\
\file{source/} & The complete ordered primitive occurrence list.\\
\file{checks/} & Reconstruction, source, startup, logical and symbolic checkers.\\
\file{evidence/} & Original and portable receipts, full startup deltas and source prefix.\\
\file{proofs/} & Strong compiler, credit proof, and independent source review.\\
\file{appendices/} & Include-ready TeX for independent review and additional results.\\
\file{research/} & Portable comparison corpus, rejected-source audit, abstract parser lemmas, bounded searches.\\
\file{reference/transfer-parent/} & Unchanged inherited parent packet, including its proof texts, licenses, and exact inputs.\\
\file{provenance/} & Pinned upstream identities, inherited foundations, and attribution.\\
\bottomrule
\end{tabular}
\end{center}

The 3,728-node checker proves the stated finite
universal propositional recovery. The source
checks prove finite syntax and control identities.
The startup checker proves a finite literal run.
The symbolic checker proves arbitrary-natural
primitive contracts by sound induction rules.
The written compiler, credit, and inherited
calibration connect these levels. None of those
individual checks should be advertised as proving
all the others.

The new runtime wrapper uses the inherited proof
engine, with a compact retained interface and
new exact ST215 inputs. Its mathematical rules
are unchanged. This packet intentionally executes
the included verifier code; attribution to its
earlier authors is retained. Historical prose
inside the unchanged parent packet describes
that parent's execution policy and is not a
claim of newly authored code here.

The source snapshot is commit
\hashline{75716498437bfc7291d5900cab2983a11fd54bac}
of ProveIt. The exact K215 ordinary proof,
dominated-calculus proof, transfer compiler,
credit proof, and original source licenses
are identified in the provenance. The new
machine theorem reuses their stated arithmetic
and numerical-calculus results. No foundation
is inferred merely from successful execution
of a finite script.

Build the PDF with two runs of
\texttt{pdflatex -halt-on-error report.tex}.
The transition appendix is generated from the
dense table by the included deterministic
script. The table files themselves are the
primary machine specification.

\begin{thebibliography}{9}
\bibitem{repo}
Vladimir Reshetnikov, \emph{ProveIt: Smaller machines for Con(ZFC)},
pinned repository snapshot, 21 October 2026.
\url{https://github.com/VladimirReshetnikov/ProveIt/blob/75716498437bfc7291d5900cab2983a11fd54bac/Computability/BusyBeaver/Research/ZFC/README.md}
\bibitem{k215}
\emph{Ordinary proof of the combined 215 machine}, same snapshot.
\url{https://github.com/VladimirReshetnikov/ProveIt/blob/75716498437bfc7291d5900cab2983a11fd54bac/Computability/BusyBeaver/Research/ZFC/candidate-k215/proof.md}
\bibitem{dominated}
\emph{Dominated-rho8 alternative: proof}, same snapshot.
\url{https://github.com/VladimirReshetnikov/ProveIt/blob/75716498437bfc7291d5900cab2983a11fd54bac/Computability/BusyBeaver/Research/ZFC/candidate-k218/alternatives/dominated/proof.md}
\bibitem{transfer}
\emph{T215 Transfer}, incoming research packet,
\file{docs/incoming/T215_Transfer_2026-10-20.zip}, same snapshot.
Its frozen transfer parent is included verbatim under
\file{reference/transfer-parent/}; particularly
\file{proofs/transfer_compiler.md},
\file{proofs/protected_credit.md}, and its source-calibration provenance.
\bibitem{mp0}
\emph{MP0 Transfer 215}, incoming research packet,
\file{docs/incoming/MP0_Transfer_215.zip}, same snapshot.
The rotating-pop experiment is the target of
Appendix~\ref{sec:rotating-obstruction}.
\end{thebibliography}

\appendix
\clearpage
% Include-only section; requires amsmath, amssymb, booktabs.
\section{Independent ST215 bootstrap and calibration review}
\label{sec:st215-independent-review}

The append-zero, start-zero ST215 table passes an independent literal
blank-tape replay and a complete synthetic-AX prefix comparison.  The
result closes the dirty-tail initialization obligation for the ordinary
short-transfer calculus.  The independent check does not infer the
initialization from a plausible local pattern: it compares the entire
finite set of one-cells and all nine logical registers.

\subsection{Exact finite evidence}

The literal table has 430 defined transitions, one for each symbol in
each of 215 nonhalting states.  State 254 is its only undefined target.
Starting in state zero with a blank tape and head at zero, after exactly
70,506 binary transitions it reaches the following complete frame:
\begin{center}
\begin{tabular}{ll}
\toprule
Quantity & Value\\
\midrule
Native dispatcher state & 84\\
Head position & $-2$\\
Header position and gap & $62$, gap zero\\
Source ordinal & 58, primitive $\mathsf{CI}(7)$\\
Source path & \texttt{00110010010110010}\\
Logical register order & $(N,T,W,A,D,H,X,Y,L)$\\
Logical register values & $(1,0,0,0,0,0,0,18,0)$\\
Spare blocks & 63, all zero-valued\\
One-cells on the entire tape & 100, from $-3$ through $225$\\
\bottomrule
\end{tabular}
\end{center}
The phrase ``from $-3$ through $225$'' gives the minimum and maximum,
not a contiguous run.  The full 100-cell set appears in the JSON receipt.

Here is an independent description of the reconstructed frame.  Relative
to the dispatcher head, place a guard at $-1$, the one-bits of the source
path at their zero-based positions, and a header at 64.  The logical-to-
physical block permutation is
\[
                         (6,2,0,4,3,7,8,5,1).
\]
Beginning at position 65, each block is a separator zero followed by
$v+1$ one-cells for its value $v$.  Encode the nine permuted logical
registers followed by the 63 zero spare blocks, and translate by $-2$.
This construction gives exactly the tape obtained by direct execution;
there are no unexamined one-cells outside the frame.

Two independently formulated source-control resolvers, one using subtree
interval endpoints and one using arithmetic cuts of binary paths, agree
on all 4,732 primitive leaves.  The resolved graph is also identical to
the complete unrotated context-root short-transfer graph.  In particular,
the AX/MP continuation at ordinal 4696, $\mathsf{CC}(3)$, returns to
ordinal zero.  This is the continuation that the archived scratch-biased
experiment failed to preserve.  Independent subtree interning counts
201 distinct source DAG nodes.

The source trace starts at AX entry ordinal 25 with
\[
                      (1,0,2,0,0,0,0,0,0).
\]
After exactly 177 source transitions it reaches ordinal 58 with
\[
                      (1,0,0,0,0,0,0,18,0),
\]
exactly the register store of the binary frame.  The complete 178-row
trace, including both endpoints, is retained in the receipt.  This
synthetic prefix is a continuation witness: the argument does not claim
that the actual blank run previously passed through that AX-entry store.
Determinism identifies their futures once they reach the same ordinal
and complete logical store, and the whole-tape simulation identifies
that source future with the actual native future.

The independent table hash agrees with the complete runtime receipt.
That receipt reports 6,151 guarded cases, all 430 literal rows covered,
1,617,630 symbolic rules, no failures, and minimum duration five.  The
independent finite checker here does not rerun those symbolic cases;
its role is to establish the exact initial frame and source-prefix bridge
to which that all-natural simulation theorem applies.

\subsection{Completion of the partial AX}

Write $P$ for Cantor pairing, $[x;w]=P(x,w)$,
$I(x,y)=4P(x,y)$, and $Q_i x=4P(i,x)+2$.  Truth has code zero and
falsity has code one.  For a command store use $(t,w,d,h)$ for top,
encoded stack, additive counter and cached formula.

The synthetic AX has parameters
\[
       (t,w,d,h)=(0,2,0,0),\qquad (c,x,y)=(0,0,0).
\]
Its inherited all-natural arithmetic contract terminates and yields
\[
  (t,d,h)=(0,0,0),\qquad N=1,\quad L=0,
\]
with stack
\[
 [r_1(0,0,0),\ldots,r_7(0,0,0),
   L(0,0,0,0),Q_0 0,0;2].
\]
The first eight displayed heads are literal axiom instances, the ninth
is the generalization $Q_0 0=2$ of truth, and the tenth is saved truth.
Thus all ten explicit heads are theorems of the calibrated numerical
calculus.  The gate is zero because $c=A=B=0$.  The remaining tail is
raw integer two; it is not assumed safe.  Indeed
$2=P(0,1)$ and $1=P(1,0)$, so that tail contains a falsity head.

Only the finite 177-step prefix is executed in the bridge check.  The
termination and result of the much larger remaining AX are supplied by
its arithmetic contract.  Literal simulation of the enormous encoded
output is unnecessary and is not claimed.

\subsection{Protected-tail soundness}

The actual source implements the following short transfer:
\[
 \mathsf S(t,[x;v],d,h)=
 \begin{cases}
    \text{halt},&h=1,\\
    (h,v,d,x),&h\ne1.
 \end{cases}
\]
It tests only the incoming cache.  In particular, it is distinct from
the longer composite transfer that also tests the incoming top.  MP and
AX retain their inherited equations.  The field arities are one for
S (tag zero), two for MP (tag one), and three for AX (tags at least two).

The all-natural scheduler credit proof applies to exactly those tags and
arities.  Give S weight $-1$, MP weight zero, and AX weight ten, ignoring
the additional zero head of a successful MP.  Starting with fresh
candidate one and an empty proof pointer, every finite schedule prefix
has weight at least $-2$.  The exceptional initial schedule is
$W(1);W(2)=[\mathsf{MP},\mathsf S,\mathsf S]$; every later complete
Cantor diagonal is prefix-safe.  Thus the ten theorem heads produced
above leave at least eight protected explicit theorem heads forever.

For clarity, the credit proof uses the exact recurrences
\[
\begin{aligned}
 W(n)&=\mathsf S;R(y) &&\text{if }\operatorname{unpair}(n)=(0,y),\\
 W(n)&=\mathsf{MP};R(f(y)) &&\text{if }\operatorname{unpair}(n)=(1,y),\\
 W(n)&=\mathsf{AX};R(f(f(y))) &&\text{if }\operatorname{unpair}(n)=(a,y),\ a\ge2,
\end{aligned}
\]
where $f$ is right projection, $R(0)$ is empty, and $R(p)=W(p-1)$
for $p>0$.  In particular, AX takes two further projections after its
selector, not three.  Its proof is an induction over complete Cantor
diagonals, not a bound inferred from a finite simulation.

Maintain a finite explicit list of theorem heads over the fixed tail
two.  AX adds ten such heads.  MP preserves the list or adds theorem
zero.  Each S has at least one explicit head available before popping,
because the conservative count after its decrement remains at least
eight.  Thus it cannot expose the tail.  Its new top and cache are
theorems; a halt witnesses theorem one.  This argument counts explicit
wrappers, not numerical code size, so it includes the alias $P(0,0)=0$.
No assertion that all raw integers are formulas, or that the buried
tail consists of formulas, is used.

\subsection{Refutational completeness}

The inherited strong compiler is already proved for this exact short
S, rather than by treating it as equivalent to old NIP.  Put
$D=\mathsf{AX}(0;0,0)$ and
\[
                    H=D;\mathsf S^{11}.
\]
It satisfies, from every natural store,
\[
 H(t,[x;w],d,h)=
 \begin{cases}
   \text{halt},&t=1,\\
   (t,w,0,x),&t\ne1.
 \end{cases}
\]
The first ten S tests inspect the new non-falsity rows; the eleventh
tests the saved incoming top.  For every theorem $X$, the compiler
constructs a finite block $B_X$ which, from every $(t,w,0,0)$, either
halts or returns exactly $(X,w,0,0)$.  The arbitrary-store quantifier
is essential here: it covers any preserved stack tail and every top
produced by a resetting AX.

If the calculus proves one, the fixed word $B_1;H$ therefore halts
from every such reset store.  Let $F(\vec a,r)$ denote right-nested
pairing of a finite field list over residual $r$, and define
\[
 \operatorname{tail}([])=0,\qquad
 \operatorname{tail}(k::V)=1+
       F(\operatorname{fields}(k),\operatorname{tail}(V)).
\]
The positive residual is decremented before decoding its command, so
this is the exact continuation encoding.  Put
$R=\operatorname{tail}(B_1;H)$.  For every $Y\in\mathbb N$, the fresh
candidate
\[
                  N_Y=P(2,P(Y,P(1,R)))
\]
begins with $\mathsf{AX}(0;1,Y)$ and then executes $B_1;H$.
The initial AX resets the counter and cache, and the strong compiler
handles its arbitrary resulting top and stack.  Therefore every $N_Y$
is successful from every retained store, and $N_Y\ge Y$ gives cofinality.

Every nonhalting residual command decreases the positive proof pointer:
the next residual is a right projection of its predecessor, with further
projections for MP and AX.  Each candidate therefore drains in finitely
many commands, and the unchanged next-candidate counter advances by
one.  From next candidate one the scheduler must reach a successful
$N_Y$, unless it halts sooner.  This proves refutational completeness
without any bound on the size of a proof or its code.

\subsection{Result and exact proof boundary}

Combining these arguments gives
\[
 \mathsf{ST215}\text{ halts on blank tape}
 \quad\Longleftrightarrow\quad
 \text{the inherited dominated numerical calculus proves }1.
\]
The established calibration of that unchanged calculus to ZFC then
gives, by ordinary reasoning formalizable in ZFC,
\[
 \mathrm{ZFC}\ \vdash\
 \bigl(\mathsf{ST215}\text{ does not halt on blank tape}
       \ \longleftrightarrow\ \operatorname{Con}(\mathrm{ZFC})\bigr).
\]
This conclusion uses the inherited arithmetic contracts and calculus
calibration, the all-natural protected-credit and strong-compiler proofs,
the exact finite startup/AX bridge, and the complete symbolic primitive
simulation.  Each supplies a different boundary; no finite test replaces
an unbounded theorem.  The source rotation preserves the full control
graph, and the bootstrap passes into an unchanged AX, so neither step
requires a new logical axiom schema or a claim that junk codes decode
to ordinary formulas.  This review does not claim that a new Lean or
Rocq endpoint has been compiled.

The accompanying independent files are
\begin{center}\small\ttfamily
check\_st215\_bootstrap.py\\
st215\_bootstrap\_independent.json
\end{center}
They bind the source and native bytes by SHA-256 and include the complete
finite witnesses.

% Include-only section; requires amsmath, amssymb, amsthm, booktabs.
\section{A productive-adjacency obstruction for rotating pop}
\label{sec:rotating-obstruction}

The archived packet \texttt{MP0\_Transfer\_215.zip} contains a physically
validated experimental 215-state table, \texttt{research/rotate-pop215/R.native.tm},
whose refutational completeness is explicitly left open. Its command
\[
 \mathsf R(t,[x;v],d,h)=(x,v,d,t)
\]
halts instead when $h=1$ or $x=1$. Here $[x;v]=P(x,v)$, with Cantor
pairing $P$, and $\operatorname{unpair}(0)=(0,0)$. The following argument
shows that the missing compiler cannot be repaired while keeping this
interface and the literal AX row order. It uses no consistency assumption,
no model of set theory, and no theoremhood invariant.

Put $I(a,b)=4P(a,b)$, $Q_a b=4P(a,b)+2$, $E(a,b)=4P(a,b)+3$,
$M(a,b)=4P(a,b)+1$, $F=1$, and $NX=I(X,F)$. Define
$\exists_aX=NQ_aNX$. The actual AX output is
\[
 (G,[r_1,r_2,r_3,r_4,r_5,r_6,r_7,L,Q_Uh,t;w],0,0),
\]
where $U=d$, $A=d+x$, $B=h+y$, $C=c$, the $r_j$ are exactly
\begin{align*}
 r_1&=\exists_AE(A,B),\\
 r_2&=I(Q_A I(B,C),I(Q_AB,Q_AC)),\\
 r_3&=I(E(A,B),I(M(C,A),M(C,B))),\\
 r_4&=I(Q_AQ_BC,Q_BQ_AC),\\
 r_5&=I(E(C,A),I(M(C,B),M(A,B))),\\
 r_6&=I(\exists_AQ_AC,C),\\
 r_7&=I(E(C,A),I(E(C,B),Q_{M(A,B)}E(A,B))),\\
 L&=I(I(I(C,U),A),I(I(A,C),I(B,C))).
\end{align*}
The gate $G$ is either zero or the literal full bundle
\[
 G=I(I(V,J),F),\qquad V=I(M(A,B),Q_CM(A,B)).
\]
The contents of $J$ are immaterial to the argument: the full inherited
term is retained. Thus the result does not weaken the underlying theory.

Call a pair $(m,a)$ \emph{productive} when
\[
             \exists b>0\quad m=I(a,b).
\]
For a command store, consider the ordered spine
$h,t,x_1,x_2,\ldots$, where $x_i$ are successive heads of its encoded
stack. A spine is \emph{blocked} when none of its entries is $1$ and
none of its adjacent pairs is productive. The all-zero spine is blocked.
This definition is numerical; it does not require that an entry decode
to a conventional formula.

\paragraph{The AX adjacency lemma.}
Every newly introduced adjacent pair of the AX spine
\[
 0,G,r_1,r_2,r_3,r_4,r_5,r_6,r_7,L,Q_Uh,t,x_1,\ldots
\]
is unproductive. The existing suffix beginning with $t$ is unchanged.
For the pair beginning with zero this follows from the exact implication
alias $I(a,b)=0\Longleftrightarrow a=b=0$. For the remaining pairs,
injectivity of $I$ reduces a possible match $m=I(a,b)$ to equality
between the antecedent of $m$ and $a$. The following table gives the
contradictory outer tags; nested comparisons use injectivity once more.

\begin{center}\small
\begin{tabular}{lll}
\toprule
Pair & Required equality & Contradiction modulo $4$\\
\midrule
$(G,r_1)$, enabled gate & $V=Q_A N E(A,B)$ & $0\ne2$\\
$(r_1,r_2)$ & $Q_A N E(A,B)=r_2$ & $2\ne0$\\
$(r_2,r_3)$ & $Q_A I(B,C)=r_3$ & $2\ne0$\\
$(r_3,r_4)$ & $E(A,B)=r_4$ & $3\ne0$\\
$(r_4,r_5)$ & $Q_AQ_BC=r_5$ & $2\ne0$\\
$(r_5,r_6)$ & $E(C,A)=r_6$ & $3\ne0$\\
$(r_6,r_7)$ & $Q_A N Q_AC=E(C,A)$ & $2\ne3$\\
$(r_7,L)$ & $E(C,A)=L$ & $3\ne0$\\
$(L,Q_Uh)$ & $I(I(C,U),A)=Q_Uh$ & $0\ne2$\\
$(Q_Uh,t)$ & $Q_Uh=I(t,b)$ & $2\ne0$\\
\bottomrule
\end{tabular}
\end{center}
For a disabled gate $G=0$, the zero-pair argument applies. All generated
entries are different from $1$: the gate and eight rows have tag zero,
and $Q_Uh$ has tag two. This includes the possible alias $L=0$; no
positivity of $L$ is assumed. The inequality $U\le A$ is not needed.

\paragraph{Preservation theorem.}
Every finite word of the rotating-pop commands, MP and AX preserves a
blocked spine and never halts when begun at a blocked spine. In
particular, this holds for every finite word from $t=h=w=0$, with
arbitrary additive counter and arbitrary raw command parameters.

\emph{Proof.} AX preserves the property by the adjacency lemma.
Rotating pop maps the spine $h,t,x_1,x_2,\ldots$ to
$t,x_1,x_2,\ldots$, a suffix. Its two tested entries are not $1$.
For MP put $y=d+x$. If it succeeds, $h=I(t,y)$, so blockedness forces
$y=0$. The returned spine is $0,0,0,x_1,x_2,\ldots$. If it fails,
the returned spine is $0,0,x_1,x_2,\ldots$. Every new pair beginning
with zero is unproductive, and every new entry is zero. In either
case blockedness is preserved. Induction on the word proves the claim.
This proof also covers every numerical zero-stack alias.\hfill$\square$

\paragraph{The actual dirty-tail initialization.}
The preceding argument has a finite-prefix form. For any fixed raw tail
$J$, maintain a finite explicit head list $X$ over $J$ such that the
finite spine $[h,t]+X$ is blocked. AX adds ten explicit heads, R consumes
one, and MP adds either zero heads or one zero head. Whenever the
current explicit head list is nonempty, the same preservation proof
applies. Thus an external scheduler credit bound protects \emph{blocked}
heads just as it protects theorem-valued heads; the latter interpretation
is unnecessary here.

For the archived R table, the supplied exact blank startup reaches the
interior of the unchanged AX at source ordinal $55$, with registers
$[1,11,2,0,0,0,0,0,0]$. Its recorded AX-prefix match completes to next
candidate $1$, empty proof pointer, $t=d=h=0$, and
\[
 [r_1(0,0,0),\ldots,r_7(0,0,0),L(0,0,0,0),Q_0 0,0;4].
\]
The displayed ten-head prefix is blocked by the lemma, independently
of the tail $4=P(1,1)$. The unchanged wire arities are one field for
R (tag zero), two fields for MP (tag one), and three fields for AX
(tags at least two). The inherited all-natural Cantor-diagonal credit
proof assigns weights $-1,0,10$ and proves that every schedule prefix
from candidate $1$ has weight at least $-2$. Consequently at least
eight explicit blocked heads remain throughout the actual execution.
The buried $1$ heads are never reached. The source therefore never
halts, by an ordinary proof formalizable in ZFC. Combining the supplied
positive-duration primitive simulation and the complete blank prefix
gives the same unconditional nonhalting statement for the exact R table.

This last exact-table conclusion is conditional on those independently
specified physical certificates; the new proof here establishes the
command-level obstruction and its protected-prefix transport. It does
not rerun the large physical certificates. The checker accompanying this
section verifies all ten literal adjacency obstructions and the exact
numerical gate/row formulas used in this proof. It does not claim to
be a proof-assistant kernel check.

\paragraph{A fixed theorem defeats the proposed strong compiler.}
The calibrated numerical calculus proves
\[
                   X=I(0,Q_0 0)=I(0,2)=20.
\]
Indeed truth gives zero, GEN gives $Q_0 0=2$, and the derived
propositional K schema gives $2\to(0\to2)$; one MP gives $X$.
The dominated calculus inherits that K schema from its established
propositional recovery certificate. A strong block for $X$, applied
to the all-zero store, would have to halt or return
$(20,0,0,0)$. The preservation theorem excludes halting. The alleged
return has spine $0,20,0,0,\ldots$, containing the productive adjacent
pair $(20,0)$ because $20=I(0,2)$ and $2>0$. Hence the return is
also impossible. No strong compiler for all numerical theorems exists,
even for this fixed elementary theorem; this conclusion does not use
any consistency or incompleteness assumption.

\paragraph{Consequence for refutational completeness.}
More generally, no block can take a blocked all-zero store to a
return state with top $1$ or to a halt. If a
refutation-complete compiler for the calibrated ZFC calculus existed
and were provable in ZFC, the unconditional nonhalting proof would
yield $\operatorname{Con}(\mathrm{ZFC})$ in ZFC. Under the usual
consistency assumption this contradicts the second incompleteness
theorem. The obstruction is specifically caused by preserving the
literal AX order while replacing transfer-pop with rotating pop.
Changing that order, emitting an additional bridge pair, or changing
the return map would need a new analysis.

% Include-only section; requires amsmath, amssymb, booktabs.
\section{A control-flow error in the archived scratch-biased experiment}
\label{sec:scratch-correction}

The scratch-biased 215-state experiment in
\texttt{T215\_Transfer\_2026-10-20.zip} has a successful quantified
primitive-runtime certificate, but its README's high-level command
description is incorrect. Replacing a continuation leaf by a two-leaf
subtree changed the continuation's target. This defect is visible in
the literal source independently of physical allocation.

The pre-insertion continuation was $\mathsf{CC}(3)$ at binary path
$10011$. Its three right ancestors are at indices $0,3,4$, so the
third outward right ancestor is the root and the continuation returns
to source ordinal zero. The replacement subtree
$(\mathsf{CI}(6),\mathsf{CC}(3))$ places the continuation at $100111$.
There are now four right ancestors, at indices $0,3,4,5$. Its third
outward right ancestor has prefix $100$, whose first leaf is source
ordinal $23$. Thus the literal edge is
\[
       4697:\ \mathsf{CC}(3)\longrightarrow23,
\]
not an edge to zero. The added right branch must be counted in a
relative continuation instruction.

\paragraph{Exact unconditional halting suffix.}
At ordinal $4696$, assume only that logical registers
$A=H=X=0$; write the arbitrary current top as $T=t$. Other logical
registers are irrelevant to the following suffix. The actual execution
is:
\begin{enumerate}
 \item Ordinal $4696$ increments scratch $X$ to one; ordinal $4697$
 returns to $23$.
 \item The zero selector $A$ takes the zero branch at $23$, reaching
 ordinal $4698$, the top-clear loop. That loop clears $T$ in exactly
 $2t+1$ source steps and does not change $H$ or $X$.
 \item The copy-cache code begins at $4700$. Its $H$ loop is skipped
 because $H=0$, while its scratch-restoration loop consumes the
 outstanding $X=1$ and increments $H$ to one.
 \item The inherited NIP test begins at $4707$ and recognizes
 precisely this $H=1$. It reaches the halt primitive at ordinal
 $4718$.
\end{enumerate}
Counting transitions from $4696$ to the halt primitive gives exactly
\[
                            2t+15.
\]
The proof is a single unary-loop induction plus the eleven fixed
post-clear edges. No inference about a very long simulated prefix is
needed.

Both AX and MP return to $4696$ with $A=H=X=0$ according to their
unchanged arithmetic contracts. In particular, the supplied 69,744-step
blank trace and 172-step AX-prefix match put the exact table inside
the AX begun with registers $[1,0,3,0,0,0,0,0,0]$. Completing that AX
returns with gate/top zero and the required zero selector/cache/scratch.
The source then reaches its halt primitive after fifteen further source
transitions. Subject to the supplied complete primitive simulation, the
archived exact binary table therefore halts on blank tape
\emph{unconditionally}. It cannot provide the requested nonhalting
equivalence when ZFC is consistent.

The correction does not dispute the primitive runtime checker: that
checker correctly implements the literal source. The failed boundary
is between the literal source and the advertised high-level parser.
Consequently the proposed first-field bias, candidate steps of one or
two, and cofinality problem do not describe this exact machine. Those
ideas may remain relevant to a repaired source that actually returns
to its scheduler, but cannot calibrate these bytes.

The accompanying independent checker resolves both complete guarded
control graphs, verifies the changed edge, records the exact fifteen-step
suffix at $t=0$, and checks the affine count for $0\le t\le31$.
The finite count check corroborates the explicit unary-loop proof; it
does not replace that proof or rerun the full physical certificates.

\begingroup\small\setlength{\parskip}{4pt}
% Include-only section; requires amsmath, amssymb.
\section{Consecutive successful codes for an abstract biased parser}
\label{sec:biased-parser-synchronization}

This section proves an arithmetic completeness tool for a possible future
encoding.  It is not a description of the archived scratch-biased table,
whose source continuation is incorrect, and it is not needed to calibrate
ST215, whose candidate counter advances normally by one.  In particular,
no additional physical state bound or implementation claim follows here.

Let $P$ be Cantor pairing and $f$ its right projection.  Consider an
abstract parser with a bit $b\in\{0,1\}$, set to zero after S and to one
after MP or AX.  At a positive continuation pointer $L$, it decodes the
first field from $L-1+b$ and reads all further command fields by ordinary
right projection.  A fresh candidate is decoded without bias.  Use
\[
\begin{split}
 \operatorname{fields}(\mathsf S)&=[0],\\
 \operatorname{fields}(\mathsf{MP}(x))&=[1,x],\\
 \operatorname{fields}(\mathsf{AX}(c;x,y))&=[c+2,y,x].
\end{split}
\]
The retained arithmetic store may be arbitrary and need not be determined
by the pointer or bias.  Assume fresh-candidate numbers increase by either
one or two; the choice of increment may depend on the entire prior run.

\paragraph{A consecutive synchronization pair.}
For any $r>0$, put
\[
 a=P(2,P(1,P(0,r))),\qquad z=P(1,a),\qquad n=P(1,z).
\]
Fresh candidate $n$ parses
\[
                  \mathsf{MP}(1);\mathsf{AX}(0;0,1)
\]
and reaches pointer $r$ with bias one.  Indeed its first command leaves
pointer $a$ and sets bias one, so the next selector is decoded from $a$,
without an effective decrement.  The three AX fields are $[2,1,0]$.

For every $z\ge0$, the elementary Cantor identity is
\[
                         P(1,z)+1=P(0,z+1).
\]
Consequently fresh candidate $n+1$ begins with S, leaving pointer $z+1$
and bias zero.  Unless that S already halts, the next selector is read
from $z=P(1,a)$, so it is MP with argument two and remaining pointer
$P(1,P(0,r))$.  Its bias is one, and the next MP therefore has argument
zero and leaves pointer $r$.  Thus the second word is
\[
                       \mathsf S;\mathsf{MP}(2);\mathsf{MP}(0),
\]
and it too reaches pointer $r$, bias one, or halts sooner.

This synchronizes only the parser position and bias.  It does \emph{not}
assert equality of the two retained arithmetic stores.  A common suffix
must therefore begin with a reset before using a compiler whose input
requires zero counter and cache.

\paragraph{Positive-tail word encoding.}
Let $F(\vec a,r)$ be right-nested pairing of a finite field list over
residual $r$, and let $\beta(k)$ be the bias produced by command $k$.
For a finite command word $V$ and a positive unused tail $q$, define
\[
\begin{split}
 E_b([],q)&=q,\\
 E_b(k::V,q)&=1-b+
       F(\operatorname{fields}(k),E_{\beta(k)}(V,q)).
\end{split}
\]
Pairing dominates its right argument, so induction gives
$E_b(V,q)\ge q>0$.  From this pointer and bias $b$, the parser subtracts
$1-b$ and obtains exactly the nested fields of $k$.  It then reaches
the suffix pointer with precisely the bias required by the induction
hypothesis.  Thus it executes all of $V$, unless a command halts earlier,
and leaves $q$ if the whole word returns.

The positivity of the unused tail is necessary.  With base tail zero,
one would obtain $E_1([\mathsf S],0)=0$, which the scheduler would treat
as a fresh fetch; the intended final S would never execute.  The positive
tail avoids this zero alias without changing any command field.

\paragraph{Cofinal successful pairs.}
Suppose a strong refutation block $B_1;H$ halts from every
$(t,w,0,0)$, as in the inherited short-transfer compiler.  Let $V$ be
an AX reset followed by that block.  It is successful from every retained
store.  For each $q\ge1$, take $r=E_1(V,q)$ and construct the consecutive
pair $n(r),n(r)+1$ above.  The first synchronization word reaches $V$
from some retained store; the second either halts or does the same.
Therefore both codes are successful from every initial retained store.
Moreover
\[
                           n(r)\ge r\ge q,
\]
so these successful consecutive pairs are cofinal.

Every strictly increasing integer sequence with gaps at most two meets
every sufficiently late consecutive pair.  Otherwise, on first crossing
$n$, it would jump from a value at most $n-1$ to a value at least $n+2$,
a gap of at least three.  This argument does not assume that the gap
choices are independent of the candidate, parser bias or arithmetic data.
Thus skipping at most one candidate cannot defeat the successful pairs.

Each nonhalting candidate also drains.  For $L>0$, the first unpairing
input $v=L-1+b$ satisfies $0\le v\le L$.  If $v=0$, its right projection
is zero and is less than $L$.  If $v>0$, then $f(v)<v\le L$; equality
$f(v)=v$ would force the triangular part of $P$ to vanish and hence
$v=0$.  Further command fields only decrease the residual.  Every step
therefore strictly decreases the positive continuation pointer.  The
scheduler reaches each subsequent selected candidate in finitely many
commands unless it has already halted.

These are all-natural arithmetic arguments.  The companion checker
\texttt{check\_biased\_sync.py} corroborates the two
parsed prefixes for 1,000 positive residuals and checks 18 positive-tail
word encodings.  Those finite checks do not establish an implementation,
protected-tail soundness, or a physical state count.  An implementation
using this lemma would still need those separate obligations.

\endgroup
\clearpage
\section{Complete dense transition table}
\label{sec:table}

Start is $0$ on the all-zero tape, and halt is
$215$. Each entry is \emph{write, direction, target}.
There is no row for halt. The numbering is the
sorted native working-state numbering, starting
at zero. The full bijection is
\file{machine/canonicalization.json}.
\begingroup\small\ttfamily
\setlength{\tabcolsep}{9pt}
\begin{longtable}{rll@{\hspace{12mm}}rll}
\toprule
q & read 0 & read 1 & q & read 0 & read 1\\\midrule
\endfirsthead
\toprule
q & read 0 & read 1 & q & read 0 & read 1\\\midrule
\endhead
\bottomrule\endfoot
0 & 1,R,91 & 1,R,0 & 108 & 0,R,148 & 1,R,183\\
1 & 0,L,1 & 1,L,2 & 109 & 0,L,204 & 0,L,94\\
2 & 0,L,3 & 1,L,2 & 110 & 0,R,123 & 1,R,103\\
3 & 0,L,5 & 1,L,2 & 111 & 0,R,206 & 1,R,180\\
4 & 1,R,6 & 0,L,4 & 112 & 0,R,162 & 1,R,178\\
5 & 0,L,9 & 1,L,2 & 113 & 0,R,146 & 1,R,144\\
6 & 0,R,6 & 0,R,10 & 114 & 0,R,157 & 1,R,197\\
7 & 0,R,83 & 1,R,187 & 115 & 0,L,130 & 0,L,1\\
8 & 0,R,14 & 1,R,13 & 116 & 0,R,38 & 1,R,72\\
9 & 0,L,15 & 1,L,2 & 117 & 0,R,61 & 1,R,131\\
10 & 0,L,16 & 0,R,10 & 118 & 0,R,61 & 1,R,132\\
11 & 0,L,26 & 1,R,128 & 119 & 0,R,70 & 0,L,16\\
12 & 0,R,19 & 1,R,20 & 120 & 0,R,14 & 1,R,131\\
13 & 0,L,88 & 1,L,31 & 121 & 0,R,21 & 1,R,133\\
14 & 0,L,11 & 0,R,22 & 122 & 0,R,145 & 1,R,175\\
15 & 0,L,23 & 1,L,2 & 123 & 0,R,135 & 1,L,20\\
16 & 1,L,1 & 0,L,1 & 124 & 0,R,129 & 1,R,177\\
17 & 0,R,24 & 0,R,32 & 125 & 0,R,201 & 1,R,168\\
18 & 0,R,198 & 1,R,197 & 126 & 0,R,75 & 1,R,173\\
19 & 0,R,28 & 1,R,29 & 127 & 0,R,17 & 1,R,212\\
20 & 0,R,30 & 1,R,31 & 128 & 0,R,75 & 1,R,134\\
21 & 0,R,24 & 0,R,59 & 129 & 0,R,12 & 1,R,139\\
22 & 0,R,33 & 1,R,22 & 130 & 0,R,30 & 1,R,69\\
23 & 0,L,34 & 1,L,2 & 131 & 0,R,38 & 1,L,123\\
24 & 0,L,35 & 0,R,0 & 132 & 0,R,151 & 0,L,16\\
25 & 0,R,36 & 1,R,25 & 133 & 0,R,152 & 0,L,16\\
26 & 0,R,127 & 1,R,37 & 134 & 0,R,158 & 1,R,137\\
27 & 0,R,38 & 0,L,39 & 135 & 0,R,24 & 0,R,47\\
28 & 0,R,40 & 1,R,41 & 136 & 0,R,138 & 1,R,177\\
29 & 0,R,42 & 1,R,43 & 137 & 0,R,191 & 1,R,93\\
30 & 0,R,44 & 1,R,45 & 138 & 0,R,95 & 1,R,156\\
31 & 0,R,21 & 1,R,46 & 139 & 0,R,172 & 1,R,163\\
32 & 0,R,47 & 1,R,32 & 140 & 0,R,12 & 1,R,156\\
33 & 0,R,32 & 1,R,33 & 141 & 0,R,87 & 1,L,163\\
34 & 0,L,49 & 1,L,2 & 142 & 0,R,87 & 1,R,155\\
35 & 1,R,35 & 1,L,1 & 143 & 0,L,150 & 1,L,160\\
36 & 0,R,50 & 1,R,36 & 144 & 0,R,83 & 1,R,87\\
37 & 0,R,13 & 1,R,52 & 145 & 0,R,156 & 1,R,12\\
38 & 0,L,11 & 0,R,36 & 146 & 0,R,82 & 1,R,141\\
39 & 0,L,4 & 1,L,53 & 147 & 0,R,161 & 1,R,165\\
40 & 0,R,24 & 0,R,25 & 148 & 0,L,66 & 1,L,139\\
41 & 0,R,13 & 0,L,16 & 149 & 0,R,146 & 1,R,199\\
42 & 0,R,55 & 1,R,56 & 150 & 0,R,156 & 1,R,113\\
43 & 0,R,57 & 1,R,58 & 151 & 0,R,96 & 1,R,96\\
44 & 0,R,24 & 0,R,54 & 152 & 0,R,48 & 0,R,32\\
45 & 0,R,60 & 1,R,61 & 153 & 0,R,68 & 1,R,200\\
46 & 0,R,62 & 0,L,16 & 154 & 0,L,207 & 1,R,95\\
47 & 0,R,25 & 1,R,47 & 155 & 0,R,165 & 1,R,176\\
48 & 0,L,35 & 0,R,63 & 156 & 0,R,19 & 1,R,167\\
49 & 0,L,74 & 1,L,2 & 157 & 0,R,170 & 1,R,162\\
50 & 0,R,16 & 1,R,50 & 158 & 0,R,156 & 1,R,95\\
51 & 0,R,145 & 1,R,141 & 159 & 0,R,124 & 1,R,161\\
52 & 0,R,67 & 1,R,68 & 160 & 0,R,96 & 1,R,159\\
53 & 0,L,39 & 1,L,53 & 161 & 0,L,7 & 1,R,140\\
54 & 0,R,22 & 1,R,54 & 162 & 0,L,108 & 1,R,170\\
55 & 0,L,11 & 0,R,50 & 163 & 0,R,144 & 1,R,169\\
56 & 0,R,60 & 0,L,16 & 164 & 0,R,193 & 1,R,28\\
57 & 0,R,69 & 1,R,70 & 165 & 0,R,158 & 1,R,141\\
58 & 0,R,61 & 1,R,71 & 166 & 0,R,175 & 1,R,175\\
59 & 0,R,54 & 1,R,59 & 167 & 0,R,178 & 1,R,121\\
60 & 0,R,48 & 0,R,25 & 168 & 0,R,143 & 1,R,162\\
61 & 0,L,28 & 1,L,72 & 169 & 0,R,75 & 1,R,60\\
62 & 0,L,116 & 1,L,30 & 170 & 0,R,143 & 1,R,87\\
63 & 0,R,73 & 1,R,63 & 171 & 0,R,103 & 1,R,95\\
64 & 0,R,96 & 1,R,148 & 172 & 0,R,138 & 1,R,83\\
65 & 1,R,181 & 1,L,53 & 173 & 0,R,172 & 1,R,75\\
66 & 0,R,142 & 1,R,143 & 174 & 0,R,20 & 1,R,37\\
67 & 0,R,76 & 1,R,215 & 175 & 0,R,140 & 1,L,147\\
68 & 0,L,193 & 1,R,31 & 176 & 0,R,96 & 1,R,51\\
69 & 0,R,40 & 1,R,62 & 177 & 0,R,165 & 1,R,51\\
70 & 0,L,88 & 1,R,77 & 178 & 0,R,17 & 1,L,130\\
71 & 0,R,55 & 1,L,78 & 179 & 0,R,76 & 1,R,115\\
72 & 0,R,79 & 1,R,80 & 180 & 0,R,105 & 1,R,123\\
73 & 1,R,0 & 1,R,63 & 181 & 0,L,39 & 1,L,109\\
74 & 0,R,74 & 1,R,81 & 182 & 0,R,168 & 1,R,156\\
75 & 0,L,143 & 1,L,139 & 183 & 0,R,148 & 1,R,114\\
76 & 0,R,84 & 1,R,85 & 184 & 0,R,11 & 1,R,97\\
77 & 0,R,88 & 1,L,131 & 185 & 0,R,189 & 1,R,98\\
78 & 0,L,107 & 1,R,56 & 186 & 0,R,144 & 1,R,195\\
79 & 0,R,48 & 0,R,33 & 187 & 0,R,154 & 1,R,188\\
80 & 0,R,90 & 0,L,41 & 188 & 0,R,172 & 1,R,182\\
81 & 0,R,194 & 1,R,174 & 189 & 0,R,64 & 1,R,7\\
82 & 0,R,95 & 1,R,12 & 190 & 0,R,148 & 1,R,126\\
83 & 0,R,87 & 1,R,96 & 191 & 0,R,153 & 1,R,111\\
84 & 0,R,203 & 1,R,21 & 192 & 0,R,150 & 1,R,122\\
85 & 0,R,115 & 1,R,99 & 193 & 0,R,89 & 1,R,87\\
86 & 0,R,148 & 1,R,112 & 194 & 0,R,214 & 1,R,115\\
87 & 0,R,101 & 1,R,102 & 195 & 0,R,202 & 1,R,108\\
88 & 0,R,79 & 1,R,78 & 196 & 0,R,148 & 1,R,166\\
89 & 0,R,100 & 1,R,171 & 197 & 0,R,196 & 1,R,157\\
90 & 0,R,104 & 1,R,105 & 198 & 0,R,138 & 1,R,136\\
91 & 0,R,65 & 0,R,0 & 199 & 0,R,192 & 1,R,18\\
92 & 0,R,27 & 1,R,106 & 200 & 0,R,87 & 1,R,205\\
93 & 0,R,179 & 0,L,204 & 201 & 0,R,160 & 1,R,190\\
94 & 1,L,109 & 1,L,94 & 202 & 0,R,125 & 1,R,149\\
95 & 0,R,19 & 1,R,110 & 203 & 0,R,44 & 1,R,99\\
96 & 0,R,60 & 1,R,60 & 204 & 0,L,4 & 0,L,204\\
97 & 0,R,164 & 1,R,154 & 205 & 0,R,193 & 1,R,138\\
98 & 0,R,66 & 1,R,186 & 206 & 0,R,161 & 1,R,207\\
99 & 0,R,104 & 0,L,204 & 207 & 0,R,208 & 1,R,179\\
100 & 0,R,116 & 1,R,117 & 208 & 0,R,141 & 1,R,209\\
101 & 0,R,71 & 1,R,58 & 209 & 0,R,142 & 1,R,210\\
102 & 0,R,57 & 1,R,118 & 210 & 0,R,211 & 1,R,86\\
103 & 0,R,21 & 1,R,119 & 211 & 0,R,185 & 1,R,162\\
104 & 0,R,21 & 0,L,10 & 212 & 0,R,184 & 1,L,72\\
105 & 0,R,61 & 1,R,41 & 213 & 0,R,92 & 1,R,107\\
106 & 0,R,120 & 1,R,8 & 214 & 0,R,213 & 1,R,127\\
107 & 0,R,100 & 1,R,121 &  & & \\
\end{longtable}
\endgroup

\end{document}
